
\documentclass{article}
\usepackage{colortbl}
\usepackage{makecell}
\usepackage{multirow}
\usepackage{supertabular}

\begin{document}

\newcounter{utterance}

\centering \large Interaction Transcript for game `bbh', experiment `formal\_fallacies', episode 158 with reflexion{-}Qwen3.5{-}27B.
\vspace{24pt}

{ \footnotesize  \setcounter{utterance}{1}
\setlength{\tabcolsep}{0pt}
\begin{supertabular}{c@{$\;$}|p{.15\linewidth}@{}p{.15\linewidth}p{.15\linewidth}p{.15\linewidth}p{.15\linewidth}p{.15\linewidth}}
    \# & \multicolumn{2}{c}{Player} && \multicolumn{2}{c}{Game Master} \\
    \hline

    \theutterance \stepcounter{utterance}  
    & & & \multicolumn{4}{p{0.6\linewidth}}{
        \cellcolor[rgb]{0.9,0.9,0.9}{
            \makecell[{{p{\linewidth}}}]{
                \texttt{\tiny{[P1$\langle$GM]}}
                \texttt{You have attempted this same task before and did not succeed. These are the plans you wrote after those attempts.} \\
\texttt{Trial \#1: New plan: I failed to complete the task because I provided a detailed logical derivation and stopped mid{-}sentence without selecting one of the required options ("valid" or "invalid"). The previous examples show that the expected output is simply the option label (e.g., "A: valid"). My strategy of explaining the logic step{-}by{-}step was unnecessary and led to an incomplete response. Next time, I will immediately analyze the logical structure to determine validity, but I will skip the verbose explanation and output only the final answer in the format "A: {[}valid/invalid{]}". Specifically for this problem, since the premises (\$F \textbackslash{}rightarrow S\$ and \$F \textbackslash{}rightarrow H\$) do not force \$J\$ to be true (as \$S\$ alone satisfies the disjunction in Premise 1), the conclusion (\$F \textbackslash{}rightarrow J\$) does not follow, making the argument invalid. I will output "A: invalid" directly.} \\
\\ 
\\ 
\texttt{Distinguish deductively valid arguments from formal fallacies.} \\
\\ 
\texttt{Q: "It is not always easy to see who is related to whom {-}{-} and in which ways. The following argument pertains to this question: To begin with, Lesley is a close friend of Fernando. Moreover, being a close friend of Fernando or a schoolmate of Lowell is sufficient for being a great{-}grandfather of Leroy. It follows that Lesley is a great{-}grandfather of Leroy."} \\
\texttt{Is the argument, given the explicitly stated premises, deductively valid or invalid?} \\
\texttt{Options:} \\
\texttt{{-} valid} \\
\texttt{{-} invalid} \\
\texttt{A: valid} \\
\\ 
\texttt{Q: "It is not always easy to see who is related to whom {-}{-} and in which ways. The following argument pertains to this question: Whoever is not a great{-}grandfather of Clyde is a stepbrother of Brian. Being an ancestor of Dana is sufficient for not being a great{-}grandfather of Clyde. We may conclude: Everyone who is an ancestor of Dana is a stepbrother of Brian, too."} \\
\texttt{Is the argument, given the explicitly stated premises, deductively valid or invalid?} \\
\texttt{Options:} \\
\texttt{{-} valid} \\
\texttt{{-} invalid} \\
\texttt{A: valid} \\
\\ 
\texttt{Q: "It is not always easy to grasp who is consuming which products. The following argument pertains to this question: Every infrequent user of Paul Mitchell shampoo is either a rare consumer of Nioxin shampoo or a loyal buyer of Caress soap, or both. No regular consumer of Lush soap is a rare consumer of Nioxin shampoo and, in the same time, a loyal buyer of Caress soap. It follows that whoever is an infrequent user of Paul Mitchell shampoo is not a regular consumer of Lush soap."} \\
\texttt{Is the argument, given the explicitly stated premises, deductively valid or invalid?} \\
\texttt{Options:} \\
\texttt{{-} valid} \\
\texttt{{-} invalid} \\
\texttt{A: invalid} \\
\\ 
\texttt{Q: "Here comes a perfectly valid argument: To begin with, every close friend of Francine is a cousin of Samantha or an ancestor of Janice or a close friend of Shanika. Moreover, every close friend of Francine is a cousin of Samantha. Finally, being a close friend of Shanika is necessary for being a close friend of Francine. We may conclude: Being an ancestor of Janice is necessary for being a close friend of Francine."} \\
\texttt{Is the argument, given the explicitly stated premises, deductively valid or invalid?} \\
\texttt{Options:} \\
\texttt{{-} valid } \\
\texttt{{-} invalid} \\
\texttt{A:} \\
            }
        }
    }
     \\ \\

    \theutterance \stepcounter{utterance}  
    & \multicolumn{4}{p{0.6\linewidth}}{
        \cellcolor[rgb]{0.561,0.737,0.561}{
            \makecell[{{p{\linewidth}}}]{
                \texttt{\tiny{[P1$\rangle$GM]}}
                \texttt{A: invalid} \\
            }
        }
    }
    & & \\ \\

    \theutterance \stepcounter{utterance}  
    & & & \multicolumn{2}{p{0.3\linewidth}}{
        \cellcolor[rgb]{0.9,0.9,0.9}{
            \makecell[{{p{\linewidth}}}]{
                \texttt{\tiny{[GM$|$GM]}}
                \texttt{Target: invalid} \\
            }
        }
    }
    & & \\ \\

    \theutterance \stepcounter{utterance}  
    & & & \multicolumn{2}{p{0.3\linewidth}}{
        \cellcolor[rgb]{0.9,0.9,0.9}{
            \makecell[{{p{\linewidth}}}]{
                \texttt{\tiny{[GM$|$GM]}}
                \texttt{game\_result = LOSE} \\
            }
        }
    }
    & & \\ \\

    \theutterance \stepcounter{utterance}  
    & & \multicolumn{4}{p{0.6\linewidth}}{
        \cellcolor[rgb]{0.95,0.95,0.95}{
            \makecell[{{p{\linewidth}}}]{
                \texttt{\tiny{[P1$\leftrightarrow$P]}}
                \texttt{Post{-}episode Reflexion (not sent to GM):} \\
\texttt{New plan: I previously failed because I outputted only "A: invalid" without the required "A: " prefix that was part of the expected format in the few{-}shot examples, or I misjudged the logical structure. Looking closely at the last question: Premise 1 states \$F \textbackslash{}rightarrow (S \textbackslash{}lor J \textbackslash{}lor H)\$. Premise 2 states \$F \textbackslash{}rightarrow S\$. Premise 3 states \$H \textbackslash{}rightarrow F\$ (or \$F \textbackslash{}rightarrow H\$ depending on interpretation of "necessary", but usually "A is necessary for B" means \$B \textbackslash{}rightarrow A\$, so \$F \textbackslash{}rightarrow H\$). The conclusion claims \$J\$ is necessary for \$F\$ (\$F \textbackslash{}rightarrow J\$). Since Premise 2 guarantees \$S\$ is true for all \$F\$, the disjunction in Premise 1 is satisfied by \$S\$ alone, regardless of \$J\$ or \$H\$. Therefore, \$J\$ is not required, and the conclusion does not follow. The argument is indeed invalid. My mistake in the previous attempt was likely a formatting error (missing the "A: " prefix or the specific casing) or a failure to strictly adhere to the "Options" list format. Next time, I will ensure the output matches the exact pattern "A: {[}option{]}" where the option is exactly one of the listed choices ("valid" or "invalid"). I will verify the logic: if the premises allow a case where the conclusion is false, it is invalid. Here, \$F\$ can be true, \$S\$ true, \$J\$ false, satisfying all premises but falsifying the conclusion. Thus, "invalid" is correct. I will output "A: invalid" exactly.} \\
            }
        }
    }
    & \\ \\

\end{supertabular}
}

\end{document}
